\problemname{Picking Apples}

IOI 2015 takes place in Almaty, which roughly means ``father of apples''. Olga lives in Almaty, and has an apple orchard with two rows of trees.
Each row has $N$ trees. Each tree has a certain number of ripe apples.

Olga starts by visiting the tree in the southwest corner (the leftmost on the bottom row), and picks all its apples.
She then walks to one of the nearest trees (to the north, east, west, or south) and picks its apples.

Your task is to compute, given how many apples are on the different trees, how many apples Olga can pick in total
if she has time to pick apples from at most $K$ trees.

\section*{Constraints}
The constraints in all test cases are $1 \le N \le 15$ and $1 \le K \le 15$. The number of apples on each
tree is between $0$ and $1000$.

\begin{enumerate}
\item In test cases worth 50 points, all trees have the same number of apples.
\item In test cases worth an additional 50 points, different trees may have different numbers of apples.
\end{enumerate}

\section*{Input}
The first line contains the integers $N$ and $K$, separated by a space.

The next line contains $N$ integers -- the number of apples on the trees in the northern row, listed from the westernmost tree to the easternmost.

The third and last line also contains $N$ integers -- the number of apples on the trees in the southern row.

\section*{Output}
Your program should output a single integer -- the number of apples Olga can pick.

\section*{Explanation of Samples}
In the sample she only has time to pick apples from two trees. The tree she starts at has $6$ apples. The tree to the north has $7$ apples, while
the tree directly to the east only has $4$ apples. She can therefore pick at most $6 + 7 = 13$ apples.
